\documentclass[10pt,letterpaper,sans]{moderncv}
\moderncvstyle{banking}
\moderncvcolor{blue}

\usepackage[scale=0.81]{geometry}
\nopagenumbers{}
\setlength{\hintscolumnwidth}{2.25cm}

\name{Krishnaa}{Vadivel}
\title{Machine Learning Engineer | High-Performance Computing}
\email{srikrishnaa.careers@gmail.com}
\phone[mobile]{516-853-5663}
\social[linkedin]{sri-krishnaa-ece-phd}
\extrainfo{\href{https://krishnaa423.github.io/personal-website/}{Website} \textbar{} \href{https://github.com/krishnaa423}{GitHub: krishnaa423} \textbar{} \href{https://leetcode.com/u/krishnaa42342}{LeetCode: krishnaa42342} \textbar{} \href{https://hub.docker.com/u/krishnaa42342}{Docker Hub: krishnaa42342}}

\begin{document}
\makecvtitle

\section{Summary}
\cvitem{}{PhD researcher and engineer with experience building machine-learning models and deploying them on high-performance computing systems for large-scale materials simulation. Experience spans PyTorch, transformers, agentic ML, GPU acceleration, and distributed compute on national supercomputing clusters.}

\section{Experience}
\cventry{2021--present}{PhD Research Assistant}{Yale University, Electrical Engineering}{}{}{\begin{itemize}%
\item Developed equivariant graph neural networks in PyTorch and PyTorch Geometric, similar to DeepH and MACE, to accelerate optical-properties calculations in material systems with automatic differentiation.
\item Built distributed and accelerator-enabled ML software in Python and C++ using CUDA, ROCm, MPI, OpenMP, and HPC environments such as Frontier and Perlmutter.
\item Applied machine learning models alongside first-principles materials simulations on supercomputing clusters, connecting model development with production scientific-computing workloads.
\item Co-authored a 2025 \textit{Journal of the American Chemical Society} paper and presented APS research in 2024, 2025, and 2026, grounding ML work in real materials-science applications.
\end{itemize}}
\cventry{2019--2021}{Software and Embedded Systems Engineer}{Probe Technologies Inc. / Weatherford Wireline Services}{}{}{\begin{itemize}%
\item Built GPU-accelerated processing and visualization tools for large waveform datasets, providing practical experience with high-throughput data pipelines and performance-oriented engineering.
\item Supported production data systems with ASP-style services and Microsoft SQL-backed tooling.
\end{itemize}}

\section{Selected Projects}
\cvitem{Materials ML for Optical Properties}{Developed equivariant graph neural networks in PyTorch Geometric to accelerate optical-properties calculations and material-energy prediction.}
\cvitem{Matrix-Multiply AI Accelerator}{Built a Verilog-based matrix-multiply accelerator project for AI-oriented datapath design and hardware-aware compute fundamentals.}
\cvitem{Agentic Scientific RAG}{Built retrieval-driven tooling that reads documentation and research papers to generate structured inputs for computational science software.}
\cvitem{Distributed ML Compute}{Combined CUDA/ROCm acceleration with HPC scheduling and scalable simulation infrastructure for ML-enabled scientific workloads.}

\section{Education}
\cventry{2021--present}{PhD, Electrical Engineering}{Yale University}{}{}{}
\cventry{2023}{MPhil, Electrical Engineering}{Yale University}{}{}{}
\cventry{2023}{MS, Electrical Engineering}{Yale University}{}{}{}
\cventry{2017--2018}{MEng, Electrical and Computer Engineering}{Cornell University}{}{}{}
\cventry{2013--2017}{BS, Electrical and Computer Engineering}{Cornell University}{}{}{}

\section{Technical Skills}
\cvitem{ML}{PyTorch, PyTorch Geometric, scikit-learn, equivariant graph neural networks, transformers, automatic differentiation}
\cvitem{Programming}{Python, C++, CUDA, JavaScript, Bash}
\cvitem{Data / Compute}{NumPy, pandas, SciPy, distributed computing, CUDA, ROCm}
\cvitem{Systems}{Linux, Docker, Docker Compose, Kubernetes, gcloud}
\cvitem{Scientific}{PETSc, SLEPc, first-principles simulation, materials modeling, physics-informed ML}

\end{document}
